
%% ═══════════════════════════════════════════════════════════════════════
%%  AUTO-GENERATED by generate_experiment2_latex.py
%%  Date: 2026-03-23 20:32
%%  Source: experiment2_results/velocity_sweep_data.mat
%%  Config: 8 velocities, 10 runs/vel, σ_p=0.05, σ_q=0.05, s_max=80
%%
%%  DO NOT EDIT — re-run the script to regenerate.
%% ═══════════════════════════════════════════════════════════════════════


\subsection{Velocity Sweep: Contouring and Lag Error Analysis}
\label{sec:velocity_sweep_analysis}

To rigorously evaluate the tracking quality of both formulations across
operating regimes, we decompose the position tracking error into its
orthogonal components: the \emph{contouring error} $e_c$ (perpendicular
distance to the reference path) and the \emph{lag error} $e_\ell$
(signed projection along the path tangent).  By the Pythagorean
relationship these satisfy
\begin{equation}
    \mathrm{RMSE}_{e_p}^{2} = \mathrm{RMSE}_{e_c}^{2} + \mathrm{RMSE}_{e_\ell}^{2},
    \label{eq:decomposition}
\end{equation}
which was verified numerically for every run in our Monte~Carlo campaign.
Both components are computed in the inertial frame to ensure a fair,
frame-neutral comparison between the DQ-MPCC (which internally operates
in the body frame via $\mathrm{Log}\bigl(\mathrm{SE}(3)\bigr)$) and
the Baseline MPCC (inertial-frame errors).

The experiment sweeps 8 maximum progress velocities
$v_\theta^{\max} \in \{4, 6, 8, 10, 12, 15, 16, 18\}$~m/s
with 10 Monte~Carlo runs per velocity per controller
($\sigma_p = 0.05$~m, $\sigma_q = 0.05$~rad),
yielding 160 simulations in total over a
Lissajous path of arc length $s_{\max} = 80$~m.

%% ─────────────────────────────────────────────────────────────────────
%%  TABLE I — Full comparative metrics
%% ─────────────────────────────────────────────────────────────────────

\begin{table*}[!t]
\centering
\caption{Comparative performance across progress velocities.
$\Delta$\% is computed as $(x_{\text{DQ}} - x_{\text{Base}})/x_{\text{Base}} \times 100$;
negative values favour DQ-MPCC.  Bold marks the better value per velocity.}
\label{tab:velocity_sweep_full}
\renewcommand{\arraystretch}{1.1}
\setlength{\tabcolsep}{3pt}
\resizebox{\textwidth}{!}{%
\begin{tabular}{r|ccc|ccc|ccc|ccc|ccc}
\toprule
& \multicolumn{3}{c|}{Contouring $\mathrm{RMSE}_{e_c}$ [m]}
& \multicolumn{3}{c|}{Lag $\mathrm{RMSE}_{e_\ell}$ [m]}
& \multicolumn{3}{c|}{Position $\mathrm{RMSE}_{e_p}$ [m]}
& \multicolumn{3}{c|}{Orientation $\mathrm{RMSE}_{e_\psi}$ [rad]}
& \multicolumn{3}{c}{Lap time $t_\mathrm{lap}$ [s]} \\
$v_\theta^{\max}$
& DQ & Base & $\Delta$\%
& DQ & Base & $\Delta$\%
& DQ & Base & $\Delta$\%
& DQ & Base & $\Delta$\%
& DQ & Base & $\Delta$\% \\
\midrule
     4 & \textbf{0.2863} & 0.3287 & -12.9\% & 0.2698 & \textbf{0.2168} & +24.4\% & \textbf{0.3934} & 0.3938 & -0.1\% & \textbf{0.3808} & 0.4972 & -23.4\% & 22.29 & 24.03 & -7.3\% \\
 6 & \textbf{0.3294} & 0.3850 & -14.4\% & 0.3332 & \textbf{0.2719} & +22.5\% & \textbf{0.4685} & 0.4713 & -0.6\% & \textbf{0.5338} & 0.6767 & -21.1\% & 15.60 & 16.47 & -5.3\% \\
 8 & \textbf{0.3869} & 0.4549 & -14.9\% & 0.4059 & \textbf{0.3083} & +31.6\% & 0.5608 & \textbf{0.5498} & +2.0\% & \textbf{0.6561} & 0.8330 & -21.2\% & 12.24 & 12.71 & -3.8\% \\
10 & \textbf{0.4340} & 0.4521 & -4.0\% & 0.4705 & \textbf{0.3495} & +34.6\% & 0.6401 & \textbf{0.5715} & +12.0\% & \textbf{0.7520} & 0.8956 & -16.0\% & 10.18 & 10.10 & +0.8\% \\
12 & \textbf{0.4793} & 0.4800 & -0.1\% & 0.5335 & \textbf{0.3904} & +36.6\% & 0.7172 & \textbf{0.6188} & +15.9\% & \textbf{0.8239} & 0.9995 & -17.6\% & 8.74 & 8.81 & -0.7\% \\
15 & 0.5312 & \textbf{0.5118} & +3.8\% & 0.6404 & \textbf{0.4262} & +50.3\% & 0.8321 & \textbf{0.6660} & +24.9\% & \textbf{0.9225} & 1.0996 & -16.1\% & 7.37 & 7.42 & -0.6\% \\
16 & 0.5482 & \textbf{0.5479} & +0.1\% & 0.6764 & \textbf{0.4391} & +54.0\% & 0.8707 & \textbf{0.7021} & +24.0\% & \textbf{0.9490} & 1.1123 & -14.7\% & 7.05 & 7.05 & -0.0\% \\
18 & \textbf{0.5828} & 0.5929 & -1.7\% & 0.7504 & \textbf{0.4738} & +58.4\% & 0.9502 & \textbf{0.7590} & +25.2\% & \textbf{0.9902} & 1.1462 & -13.6\% & 6.51 & 6.49 & +0.4\% \\
\midrule
    \textit{Avg.} & 0.4473 & 0.4692 & -4.7\% & 0.5100 & 0.3595 & +41.9\% & 0.6791 & 0.5915 & +14.8\% & 0.7510 & 0.9075 & -17.2\% & 11.25 & 11.63 & -3.3\% \\
\bottomrule
\end{tabular}}
\end{table*}

%% ─────────────────────────────────────────────────────────────────────
%%  TABLE II — Orthogonal decomposition
%% ─────────────────────────────────────────────────────────────────────

\begin{table}[!t]
\centering
\caption{Orthogonal decomposition of the position error into
contouring (\%$e_c^2$) and lag (\%$e_\ell^2$) components.
The ratio $e_\ell/e_c$ indicates the dominant error type:
$>1$ implies lag-dominated tracking, $<1$ implies contouring-dominated.}
\label{tab:decomposition}
\renewcommand{\arraystretch}{1.1}
\begin{tabular}{r|ccc|ccc}
\toprule
& \multicolumn{3}{c|}{DQ-MPCC}
& \multicolumn{3}{c}{Baseline MPCC} \\
$v_\theta^{\max}$ & \%$e_c^2$ & \%$e_\ell^2$ & $e_\ell/e_c$
                  & \%$e_c^2$ & \%$e_\ell^2$ & $e_\ell/e_c$ \\
\midrule
     4 & 53.0 & 47.0 & 0.942 & 69.7 & 30.3 & 0.660 \\
 6 & 49.4 & 50.6 & 1.011 & 66.7 & 33.3 & 0.706 \\
 8 & 47.6 & 52.4 & 1.049 & 68.6 & 31.4 & 0.678 \\
10 & 46.0 & 54.0 & 1.084 & 62.6 & 37.4 & 0.773 \\
12 & 44.7 & 55.3 & 1.113 & 60.2 & 39.8 & 0.813 \\
15 & 40.8 & 59.2 & 1.206 & 59.1 & 40.9 & 0.833 \\
16 & 39.7 & 60.3 & 1.234 & 60.9 & 39.1 & 0.801 \\
18 & 37.6 & 62.4 & 1.287 & 61.0 & 39.0 & 0.799 \\
\bottomrule
\end{tabular}
\end{table}

%% ─────────────────────────────────────────────────────────────────────
%%  FIGURE — Pareto frontiers (2×2)
%% ─────────────────────────────────────────────────────────────────────

\begin{figure}[H]
\centering
\includegraphics[width=\columnwidth]{Fig/fig_pareto_quad.pdf}
\caption{Speed--accuracy Pareto frontiers for the DQ-MPCC and
Baseline MPCC across 8 progress velocities.
Each point is the median over 10 Monte~Carlo runs;
error bars denote the interquartile range (P25--P75).
(a)~Contouring RMSE, (b)~Lag RMSE, (c)~Position RMSE,
(d)~Orientation RMSE, all vs.\ lap time.
Annotations indicate $v_\theta^{\max}$ [m/s].}
\label{fig:pareto_quad}
\end{figure}

%% ─────────────────────────────────────────────────────────────────────
%%  FIGURE ANALYSIS
%% ─────────────────────────────────────────────────────────────────────

Fig.~\ref{fig:pareto_quad} presents the speed--accuracy Pareto frontiers
for both controllers.
%
In panel~(a), the DQ-MPCC (blue) lies below the Baseline (red) for $v_\theta^{\max} < 15$~m/s, after which the curves nearly merge (gap $<0.010$~m at $18$~m/s),
confirming that the coupled $\mathrm{SE}(3)$ formulation yields tighter
geometric adherence to the path.
%
Panel~(b) reveals the opposite: the Baseline's lag curve is consistently
lower, with an absolute gap that widens from
$0.053$~m at $4$~m/s
to $0.277$~m at $18$~m/s,
reflecting the decoupled formulation's ability to independently minimise
along-path delay.
%
Panel~(c) combines both components via~\eqref{eq:decomposition}:
the curves cross near $v_\theta^{\max} \approx 8$~m/s, with the lag-dominated DQ-MPCC yielding higher total position error
at high speeds.
%
Finally, panel~(d) shows that the DQ-MPCC maintains consistently lower
orientation error across the entire velocity range
($13.6$--$23.4\%$),
a direct benefit of the pose-level coupling in the cost function.

%% ─────────────────────────────────────────────────────────────────────
%%  ANALYSIS  (all numbers auto-generated)
%% ─────────────────────────────────────────────────────────────────────

\subsubsection{Contouring vs.\ Lag Decomposition}

Table~\ref{tab:velocity_sweep_full} shows that the DQ-MPCC achieves lower
contouring error for $v_\theta^{\max} < 15$~m/s
(avg.\ $-4.7\%$, best $-14.9\%$ at $8$~m/s),
whereas the Baseline MPCC dominates in lag error
(avg.\ $+41.9\%$, up to $+58.4\%$ at $18$~m/s).
The decomposition in Table~\ref{tab:decomposition} explains this asymmetry:
the Baseline is \emph{contouring-dominated}
($59$--$70\%$ of $\|e_p\|^2$ comes from $e_c$,
$e_\ell/e_c < 1$ at all speeds),
while the DQ-MPCC transitions from near parity
($53$/$\ 47$\%
at $4$~m/s) to \emph{lag-dominated}
($38$/$\ 62$\%,
$e_\ell/e_c = 1.29$ at $18$~m/s).
%
The root cause is structural: the DQ-MPCC's $\mathrm{SE}(3)$ logarithmic
error map $\boldsymbol{\rho} = \mathbf{J}_l(\boldsymbol{\varphi})^{-1}
\mathbf{R}_d^\top \Delta\mathbf{p}$ couples position and orientation,
preventing the solver from independently minimising the tangential projection.
The Baseline, by contrast, penalises $e_c$ and $e_\ell$ as decoupled cost
terms, enabling aggressive lag reduction at the expense of orientation fidelity.

